\scanpage{102}
A vrai dire, il ne serait pas difficile d'exiger une condition 5°) qui mette ensemble 3°) et 4°), quand on se donne non $X'_1\to X'_2$, mais
\diagram{\begin{tikzcd}[ampersand replacement=\&]
X'_1 \arrow[r,"\deleted{\unsure{f_1}}"] \arrow[dr,"f_1"'] \& \deleted{\unsure{Y_1}} \arrow[d,"\deleted{f}"] \\
X'_2 \arrow[r,"f_2"'] \& Y_{\unsure{2}}
\end{tikzcd}}
avec $f_2$ \deleted{\unsure{ét.}} induisant un hom \deleted{\unsure{fibrant}} surjectif sur les $\pi_1$, de sorte que $\widehat{\Pi}_1(Y,f_2(I'_2))$ peut s'interpréter comme un quotient de $\widehat{\Pi}_1(X'_2,I'_2)$, de sorte que l'action \deleted{de} $\Theta'_2$ de $\Theta$ sur ce dernier passe à $\widehat{\Pi}_1(Y,f_2(I'_2))$ \ldots

Mais je ne suis pas sûr si j'aurai besoin de cette généralisation commune de 3° et 4° (et $\unsure{=}$ de 2°, 3°, 4°, à vrai dire !).

Les conditions 2° à 4° sont pour moi \unsure{surtout|sans doute} une \uline{\unsure{recette}} commode pour trouver des conditions nécessaires sur les paramètres virtuels qui \unsure{repèrent|représentent} un $\unsure{\gamma\in\Gamma}$ indéterminé, pour qu'ils soient effectifs i.e. proviennent bel et bien d'un $\unsure{\gamma|\sigma}$. Le point de vue n'est nullement de dégager \deleted{\unsure{des}} conditions néc. et suff. qui \unsure{permettra} de \uline{caractériser} les $\Theta$ provenant de $\Gamma$ \ldots
\mspage{101}
